\documentclass[11pt,a4paper]{article}

\usepackage[utf8]{inputenc}
\usepackage[T1]{fontenc}
\usepackage[french]{babel}
\usepackage{amsmath}
\usepackage[margin=2.3cm]{geometry}
\usepackage{booktabs}
\usepackage{array}
\usepackage{ragged2e}
\usepackage{titlesec}
\usepackage{setspace}
\usepackage[colorlinks=true,linkcolor=black,urlcolor=black]{hyperref}

\titleformat{\section}{\large\bfseries}{\thesection. }{0pt}{}
\titleformat{\subsection}{\normalsize\bfseries}{\thesubsection. }{0pt}{}

\setstretch{1.15}
\newcolumntype{L}[1]{>{\RaggedRight\arraybackslash}p{#1}}

\begin{document}

% ======================= TITRE =======================
\begin{center}
  {\LARGE\bfseries NOTE METHODOLOGIQUE\par}
  \vspace{0.4em}
  {\large Analyse de la completude logique des interviews COSO\par}
  \vspace{0.3em}
  {\normalsize Base $\rightarrow$ \texttt{Questionnaire\_COSO\_VF.dta} (1032 interviews)\par}
\end{center}
\vspace{0.6em}
\hrule
\vspace{0.8em}

La presente note decrit la methode utilisee pour determiner, parmi les
\textbf{1032 interviews} de la base, celles qui sont reellement et
totalement remplies au regard de la logique conditionnelle du questionnaire.
Le detail des questions manquantes a l'echelle de chaque interview
(identifiant + question concernee) est fourni dans le fichier Excel joint :
\texttt{manquants\_id\_detail.xlsx}.

% ======================= 1. QUESTION DE DEPART =======================
\section{Question de depart}

Combien d'interviews sont completement realisees ? Dans un questionnaire a
logique de saut, certaines questions ne sont posees qu'en fonction des
reponses anterieures. Une interview est dite \emph{complete} lorsque toutes
les questions qui devaient lui etre posees --- compte tenu des reponses
precedentes --- ont effectivement recu une reponse. A l'inverse, une
interview interrompue se reconnait a des questions attendues mais laissees
vides (abandon, refus, erreur de collecte).

% ======================= 2. PRINCIPE =======================
\section{Principe de la methode}

Pour chaque interview $i$, on reconstitue l'ensemble $Q(i)$ des questions
qui devaient lui etre posees en appliquant la logique de saut du
questionnaire (sections A a R). On compte ensuite les questions attendues
qui ont reellement une reponse non vide. Le taux de completude vaut :
%
\begin{center}
  $\displaystyle
  \text{taux\_completude}(i) \;=\;
  \frac{\text{nombre de questions attendues et repondues}}{|Q(i)|}
  $
\end{center}

Deux precautions sont prises. D'une part, une reponse n'est comptee que si
la cellule contient une \emph{vraie} valeur : la base comporte des chaines
vides \texttt{""} ou \texttt{.} qui ne sont pas \texttt{NaN} et ne doivent
pas etre confondues avec une reponse. D'autre part, les champs techniques
(horodatages, coordonnees GPS), les champs d'identification et les champs
libres non soumis a la logique (C6, C9, C10, C11) sont exclus du calcul.

% ======================= 3. LOGIQUE DE SAUT =======================
\section{Logique de saut retenue}

Chaque regle est deduite du questionnaire Draft 8 et verifiee sur les
donnees. Les principales regles sont synthetisees dans le
tableau~\ref{tab:regles}.

\begin{table}[h]
\centering
\caption{Principales regles de saut du questionnaire.}
\label{tab:regles}
\small
\begin{tabular}{L{1.6cm}L{4.4cm}L{6.8cm}}
\toprule
\textbf{Section} & \textbf{Champ} & \textbf{Condition d'attente} \\
\midrule
A & A1, A4 (resultat de l'appel) & toutes les interviews \\
A & A4X (preciser le resultat)    & si A4 = autre (8) \\
A & A5 (jour / heure de rendez-vous) & si A4 = prise de rendez-vous (5) \\
C & C5 (diplome)                  & si niveau d'etudes au moins primaire (C4 $\geq$ 2) \\
C & C8 (pays)                     & si C7 = Cote d'Ivoire (1) \\
E & E2, E3, E3A                   & E2 si E1 = non ; E3 si E2 = non ; E3A si E3 = oui \\
F,G,H,I & section emploi           & si la personne a une activite (EMPLOYE = 1) \\
H & H1 a H3 (salaires)            & si salarie (F1 = 1) \\
H & H3A (classe de salaire)       & si montant refuse (H3 = 99999) \\
H & H4, H5 (revenus)              & si non salarie (F1 = 2, 3 ou 4) \\
H & H5A (classe de revenu)        & si montant refuse (H5 = 999999) \\
J & toute la section              & si sans activite (EMPLOYE = 0) \\
K & K2, K3                        & K2 si K1 = non ; K3 si K1 = oui \\
K & K4, K5, K5X                   & si K2 = oui ; K5X si K5 = autre \\
L & L4, L5 (epargne)              & si L3 = oui \\
L & L7 a L9 (dette)               & si L6 = oui \\
N & N2 a N5 (chocs)               & si N1 = oui \\
O & O2 (formation)                & si O1 = oui ; O4-O6 si O3 = oui \\
R & R1A, R2A, R3A, R3B            & si le code correspondant est choisi \\
\bottomrule
\end{tabular}
\end{table}

% ======================= 4. VALIDATION =======================
\section{Validation statistique de la methode}

Reconstruire une logique de saut comporte un risque : imaginer des regles
trop strictes (qui penaliseraient des interviews completes) ou trop
laxistes (qui masqueraient des abandons). Pour controler ce risque, les
regles sont testees sur un groupe de reference objectif : les interviews
validees par le superviseur et dont le questionnaire est rempli
(\texttt{interview\_\_status} = 100 et A4 = 1), soit \textbf{425 interviews}.
Ce groupe de reference n'entre \emph{pas} dans le calcul du taux : il sert
uniquement de metre etalon. Si la reconstruction est exacte, quasi toutes
doivent atteindre 100~\%.

\begin{table}[h]
\centering
\caption{Statistiques de validation sur les 425 references.}
\label{tab:validation}
\small
\begin{tabular}{L{8.4cm}c}
\toprule
\textbf{Indicateur de validation} & \textbf{Valeur} \\
\midrule
Interviews de reference (metre etalon)                & 425 \\
dont taux de completude $= 100\%$                     & 406  (95,5~\%) \\
mediane du taux                                       & 1,000 \\
moyenne du taux                                       & 0,997 \\
minimum / maximum du taux                             & 0,062 / 1,000 \\
references sous 90~\%                                 & 2 \\
\bottomrule
\end{tabular}
\end{table}

Un taux median de \textbf{1,000} et \textbf{95,5~\%} de references a
exactement 100~\% confirment que la logique retenue reproduit le
questionnaire reellement applique : elle n'est ni trop stricte, ni trop
laxiste.

% ======================= 5. CLASSIFICATION =======================
\section{Classification des interviews}

A partir du taux de completude, chaque interview est rangee dans l'une des
cinq classes suivantes.

\begin{table}[h]
\centering
\caption{Classes de completude et intervalles du taux.}
\label{tab:classes}
\small
\begin{tabular}{L{3.4cm}L{4.0cm}L{5.4cm}}
\toprule
\textbf{Classe} & \textbf{Intervalle du taux} & \textbf{Lecture} \\
\midrule
1-COMPLET       & taux $= 100\%$                & aucune question attendue vide \\
2-QUASI\_COMPLET & 90~\% $\leq$ taux $< 100\%$  & quasi tout est rempli, quelques manques \\
3-MOYEN         & 50~\% $\leq$ taux $< 90\%$   & une partie notable est vide \\
4-INCOMPLET     & 0~\% $<$ taux $< 50\%$       & la majorite des questions attendues sont vides \\
5-VIDE          & taux $= 0\%$                 & aucune question exploitable repondue \\
\bottomrule
\end{tabular}
\end{table}

% ======================= 6. RESULTATS =======================
\section{Resultats}

\subsection{Repartition des 1032 interviews}

Le tableau~\ref{tab:repartition} presente la repartition des cinq classes.

\begin{table}[h]
\centering
\caption{Repartition des interviews par classe de completude.}
\label{tab:repartition}
\small
\begin{tabular}{lcc}
\toprule
\textbf{Classe} & \textbf{Effectif} & \textbf{Part} \\
\midrule
1-COMPLET       & 627 & 60,8~\% \\
2-QUASI\_COMPLET & 62  & 6,0~\% \\
3-MOYEN         & 10  & 1,0~\% \\
4-INCOMPLET     & 291 & 28,2~\% \\
5-VIDE          & 42  & 4,1~\% \\
\midrule
TOTAL           & 1032 & 100,0~\% \\
\bottomrule
\end{tabular}
\end{table}

Le taux moyen global est de \textbf{0,686} et la mediane de \textbf{1,000}.
Reponse a la question de depart : \textbf{627 interviews (60,8~\%)} sont
totalement remplies, c'est-a-dire que toutes les questions attendues selon
la logique du questionnaire ont ete repondues. En y ajoutant les
quasi-completes, la base << fiable >> compte \textbf{689 interviews
(66,7~\%)}.

% ======================= 7. EXCEL =======================
\section{Exploitation du fichier Excel joint}

\paragraph{Le detail.}
Le fichier \texttt{manquants\_id\_detail.xlsx} contient une ligne par couple
(interview ; question manquante), avec l'identifiant de l'interview
(\texttt{interview\_\_id} / \texttt{interview\_\_key}), la classe de
completude et la question concernee. Il permet de :
\begin{itemize}
  \item filtrer par classe (ex. \texttt{4-INCOMPLET} ou \texttt{5-VIDE}) ;
  \item recuperer la liste des identifiants d'interviews a revoir, a
        recontacter ou a relancer ;
  \item mesurer, pour une question donnee, combien d'interviews la laissent
        vide.
\end{itemize}

\paragraph{La base.}
La base \texttt{Questionnaire\_COSO\_VF\_avec\_completude.dta} associe en
outre a chaque interview deux colonnes filtrables :
\texttt{taux\_completude} (proportion de 0 a 1) et
\texttt{classe\_completude} (de \texttt{1-COMPLET} a \texttt{5-VIDE}).

% ======================= 8. LIMITES =======================
\section{Limites et precautions}

\begin{itemize}
  \item Un champ << preciser >> facultatif (ex. \texttt{K5X}) peut manquer
        sans que le questionnaire soit reellement inacheve ; la classe
        \texttt{2-QUASI\_COMPLET} permet de recuperer ces cas.
  \item Les 425 references confirment empiriquement la validite de la
        logique, mais ne constituent pas une demonstration mathematique. La
        confiance provient de la convergence entre la logique du
        questionnaire et le comportement observe des donnees.
  \item La completude mesure que tout est repondu, pas la qualite des
        reponses : une valeur aberrante reste une valeur presente. Une etape
        complementaire serait necessaire pour verifier la coherence interne
        des reponses.
\end{itemize}

\end{document}